\documentclass[10pt,a4paper]{amsart}
\usepackage[margin=19mm]{geometry}
\usepackage{amsmath,amssymb,mathtools}
\usepackage[scr=rsfs,cal=euler]{mathalfa}
\usepackage{xfrac}
\usepackage[unicode,hidelinks,bookmarks=false]{hyperref}
\newcommand{\ie}{i.e.\ }
\newcommand{\cf}{cf.\ }
\newcommand{\resp}{respectively\ }
\newcommand{\Subsection}{\subsection}
\providecommand{\nobreakdash}{\nobreak\hbox{-}}
\setlength{\emergencystretch}{2em}
\multlinegap0pt
\usepackage{bm}
\usepackage{bbm}
\usepackage{dsfont}
\usepackage{xspace}
\usepackage{cancel}
\usepackage{mathtools}
\usepackage{amsthm}
\makeatletter
\@ifundefined{thm}{\newtheorem{thm}{Thm}}{}
\@ifundefined{lem}{\newtheorem{lem}[thm]{Lemma}}{}
\makeatother
\newcommand{\E}{\mathop{}\!\mathbb{E}}
\newcommand{\N}{\mathbb{N}}
\newcommand{\e}{\operatorname{e}} 
\newcommand{\ind}[1]{\mathbf{1}_{\{#1\}}}
\renewcommand{\P}{\mathop{}\!\mathbb{P}}
\renewcommand{\d}{\mathrm{d}}
\title[Random L\'evy Looptrees and L\'evy Maps]{Random L\'evy Looptrees and L\'evy Maps}
\author{Igor Kortchemski, Cyril Marzouk}
\date{Source: arXiv:2402.04098v3 (CC BY 4.0)}
\begin{document}
\maketitle

\noindent\emph{Source and licence.} Random L\'evy Looptrees and L\'evy Maps. Version arXiv:2402.04098v3 (CC BY 4.0), distributed under the Creative Commons Attribution 4.0 International licence, as recorded in the arXiv licence metadata for this exact version and shown on the abstract page https://arxiv.org/abs/2402.04098v3. Full rights notice: licenses/arxiv-2402.04098-CC-BY-4.0.txt.\par\smallskip

\noindent\textit{Editorial scope.} Counted unit: the Lem whose source label is \texttt{lem:ruine} in the article (source0.tex lines 1602--1611, source label \texttt{lem:ruine}), delivered together with the article's own complete original proof of it (source0.tex lines 1665--1815, one \texttt{proof} environment of 151 lines). No mathematics is written, repaired or re-derived: statement and proof are the article's own text line for line. The proof is detached from the statement in the source: the article states the result at lines 1602--1611 and prints its proof at lines 1665--1815, and the proof environment's own header names the statement, so the association is the article's own. Required context retained uncounted: lem:technique (lines 1548--1554); eq:K\_r\_lemme\_technique (lines 1597--1599). Every definition, display and label of those lines is retained unchanged. Omitted from this package and not redistributed: 1--1547, 1555--1596, 1600--1601, 1612--1664, 1816--2874. Nothing inside a retained excerpt is omitted or altered: every retained body is delivered exactly as the article's lines are written, so no omission receipt of any kind accompanies this package. Citations: the retained excerpts carry no citation command at all, so no bibliography entry of the article is transcribed and no reference is redistributed. Reference adaptation: none was needed and none was made; every reference inside the delivered unit resolves inside it, so no unresolved reference or citation is produced. Printed slips kept literal: no printed word, symbol, hypothesis or number of the counted statement or of its proof was altered. Layout and class: the article's own class is \texttt{article} with packages including inputenc, babel, amsthm, amssymb, bbm, stmaryrd, libertinus, fontenc, libertinust1math, mathalpha, cancel, microtype, soul, numprint; the wrapper is the lane's pinned \texttt{amsart} editorial wrapper at 10pt on A4, which restates the theorem-like environments the retained text needs and adds the article's own macro definitions that the retained text needs (\texttt{\textbackslash E}, \texttt{\textbackslash N}, \texttt{\textbackslash P}, \texttt{\textbackslash d}, \texttt{\textbackslash e}, \texttt{\textbackslash ind}), with extra wrapper packages (bm, bbm, dsfont, xspace, cancel, mathtools). The delivered numbering is the unit's own. Assets: no retained span contains a figure environment, a graphics-inclusion command, a PDF-page inclusion, a table or a TikZ picture; no image file is copied into this package, the package declares no asset dependency and no image file is redistributed with it. Licence: version arXiv:2402.04098v3 of this article is distributed under the Creative Commons Attribution 4.0 International licence, as recorded in the arXiv licence metadata for this exact version and shown on the abstract page https://arxiv.org/abs/2402.04098v3. A full rights notice is delivered at licenses/arxiv-2402.04098-CC-BY-4.0.txt. Characters: every retained excerpt is pure ASCII, so the delivered engine, XeLaTeX with its default Latin Modern text font, reports no missing character.
\begin{lem}\label{lem:technique}
For every $\varepsilon > 0$, almost surely, for every $r>0$ small enough, the following alternative holds: 
\begin{itemize}
\item either there exists $t \in \mathcal{C} \cap [0, r^{2(1-\varepsilon)}]$ such that $W_t \leq -r$,
\item or there exist $i \geq 1$ such that $g_i < r^{2(1-\varepsilon)}$, and $s \in [0, \min\{\delta^L_i, r^{2(1-\varepsilon)}\}]$, and also finally $t \in [0, \min\{\delta^R_i, r^{2(1-\varepsilon)}\}]$ such that both $W^{L,i}_s \leq -r$ and $W^{R,i}_t \leq -r$.
\end{itemize}
\end{lem}

\begin{equation}\label{eq:K_r_lemme_technique}
K_r \leq \frac{\log(1/r)}{(\log\log(1/r))^3}
.\end{equation}

\begin{lem}\label{lem:ruine}
Almost surely, for every $r \in \{2^{-n}, n \in \N\}$ small enough, the following assertions hold:
\begin{enumerate}
\item The path $W$ reaches $-r \log^{10 K_{r}}(1/r)$ before time $r^{2(1-1/\log\log(1/r))}$, 
\item For every $k \leq K_r$, the path $W$ reaches $-r \log^{10 k}(1/r)$ before $r \log^{10 k + 3}(1/r)$,
\item 
There exists a universal constant $c \in (0,1)$ such that for at least $c K_r$ different indices $k \leq K_r$ the following holds: 
at the first time the path $W$ reaches $-r \log^{10 k}(1/r)$, it was below $-r$ for a duration at least $(r \log^{10 k}(1/r))^2$ and it previously never reached $r \log^{10 k}(1/r)$.
\end{enumerate}
\end{lem}

\begin{proof}[Proof of Lemma~\ref{lem:technique}]
Let us work on the event that $W_t > -r$ for all $t \in \mathcal{C} \cap [0, r^{2(1-\varepsilon)}]$ and let us prove that necessarily the second case of the alternative occurs.
By monotonicity, it is sufficient to prove the claim when $r$ is restricted to the set $\{2^{-n}, n \in \N\}$. 
Let us implicitly assume also that the three assertions from Lemma~\ref{lem:ruine} hold, and let us prove that there exists $k \leq K_r$, the integer from~\eqref{eq:K_r_lemme_technique}, such that on the interval of $\mathcal{C}^c$ on which $W$ reaches $-r \log^{10k}(1/r)$ for the first time, the two bridges that we sample also reach $-r$.

 
Let us introduce some notation. For every $k \geq 1$, denote by $G_k$ and $D_k$ the boundaries $g_i$ and $d_i$ of the unique interval $(g_i, d_i) \subset \mathcal{C}^c$ such that 
$W$ reaches the level $-r \log^{10 k}(1/r)$ for the first time in the interval $(g_i, d_i)$. 
Note that these intervals are not necessarily distinct from each others and that when $r$ is small enough, by Lemma~\ref{lem:ruine}, they all intersect $(0, r^{2(1-\varepsilon)})$ for $k \leq K_{r}$, and only $D_{K_r}$ may exceed $r^{2(1-\varepsilon)}$. 
Let $\Delta_k = D_k-G_k$ and let $\Delta^L_k, \Delta^R_k \geq \Delta_k$ denote the corresponding extensions; let us change the notation and write $W^{L,k}$ and $W^{R,k}$ for the Brownian bridges with duration $\Delta^L_k$ and $\Delta^R_k$ respectively, with initial and terminal value $W_{G_k}$ and $W_{D_k}$ respectively.
We shall upper bound the probability of the event
\[A_r = \bigcap_{k \leq K_r} \Bigl(\Bigl\{\inf_{s \in [0,\Delta^L_k] \cap (0, r^{2(1-\varepsilon)})} W^{L,k}_s > -r\Bigr\} \cup \Bigl\{\inf_{t \in [0,\Delta^R_k] \cap (0, r^{2(1-\varepsilon)})} W^{R,k}_t > -r\Bigr\}\Bigr)\]
by a quantity that only depends on $r$ and whose sum over all values of $r \in \{2^{-n}, n \geq 1\}$ is finite, such as $\log^{-2}(1/r)$. 
We then conclude as in the previous proof from the Borel--Cantelli lemma. 
\medskip

\textsc{Step 1: Proof for Brownian motions.}
Let us first prove the claim when the Brownian bridges $W^L$ and $W^R$ are replaced by Brownian motions. Precisely, for every $k \leq K_r$, conditionally on $\Delta_k$, $\Delta_k^L$, $\Delta_k^R$, and $W_{G_k}$, let $X^{L,k} = (X^{L,k}_t ; t \in [0, \Delta_k^L])$ and $X^{R,k} = (X^{R,k}_t ; t \in [0, \Delta_k^L])$ be two independent Brownian motions started from $X^{L,k}_0=X^{R,k}_0=W_{G_k} > -r$.

We first claim that if the interval $[G_k,D_k]$ is too long, then the Brownian motions $X^{L,k}$ and $X^{R,k}$ on it cannot reasonably stay above $-r$. Specifically, if we denote by $X$ a generic Brownian motion independent of the rest, then by the same reasoning as in the previous proof, we have for every $k$:
\begin{align*}
&\P\Bigl(\Delta_k > 2(r+W_{G_k})^2 \log^6(1/r) \enskip\text{and}\enskip \max\Bigl\{\inf_{t \in [0,\Delta_k/2]} X^{L,k}_t, \inf_{t \in [0,\Delta_k/2]} X^{R,k}_t\Bigr\} > -r \Bigr)
\\
&\leq 2 \P\bigl(\inf\{X_s , 0 \leq s \leq (r+W_{G_k})^2 \log^6(1/r)\} > -(W_{G_k}+r) \bigr)
\\
&= 2 \P\bigl(|X_{(r+W_{G_k})^2 \log^6(1/r)}| < W_{G_k}+r \bigr)
\\
&= 2 \P(|X_1| < \log^{-3}(1/r) )
\\
&= O(\log^{-3}(1/r))
.\end{align*}
Since $K_r = o(\log(1/r))$, then a union bound combined with the Borel--Cantelli lemma imply that almost surely, all the events for $k \leq K_r$ fail for $r$ small enough (of the form $r=2^{-n}$).
Thus, in order to avoid that both $X^{L,k}$ and $X^{R,k}$ reach $-r$, we must have $\Delta_k \leq 2(r+W_{G_k})^2 \log^6(1/r)$ for every $k \leq K_r$. 
Since $G_k$ is smaller than the first hitting time of $-r \log^{10k}(1/r)$ by definition, then combined with Lemma~\ref{lem:ruine}, we have:
\[-r < W_{G_k} \leq r \log^{10k+3}(1/r)
\quad\text{and so}\quad
\Delta_k \leq 3(r \log^{10k+6}(1/r))^2 \quad\text{for each}\enskip k \leq K_r.\]

According to the strong Markov property applied at the first hitting time of $-r \log^{10k}(1/r)$, the path $W$ behaves after this time as an unconditioned Brownian motion started from this value during the remaining time of the interval $[G_k,D_k]$.
On the previous event, this remaining time is at most $3(r \log^{10k+6}(1/r))^2$.
This is not enough for $W$ to further reach $-r \log^{10(k+1)}(1/r)$. Indeed, if $X$ is a generic Brownian motion, then:
\begin{align*}
\P\Bigl(\inf_{[0,4(r \log^{10k+6}(1/r))^2]} X < -r \log^{10k+8}(1/r)\Bigr)
&= \P\Bigl(|X_1| > \frac{1}{2} \log^{2}(1/r)\Bigr)
\\
&\leq 2 \e^{- \frac{1}{4} \log^{2}(1/r)}
,\end{align*}
and the right-hand side, multiplied by $K_r$, is again a convergent series when summing over all values $r \in \{2^{-n}, n \in \N\}$.
Hence, we may also assume that the intervals for different values of $k \leq K_r$ are disjoint, and so the Brownian motions $X^{L,k}$ and $X^{R,k}$ that we sample are independent.

Finally, by the last claim of Lemma~\ref{lem:ruine}, for a positive proportion of the indices $k \leq K_r$ we have both $W_{G_k} \leq r \log^{10k}(1/r)$ and $\Delta_k > (r \log^{10k}(1/r))^2$, so there is a probability uniformly bounded away from $0$ that both $X^{L,k}$ and $X^{R,k}$ reach $-r$, in the first half of the interval.
We conclude by independence that this occurs for a positive proportion of the indices $k \leq K_r$.
\medskip

\textsc{Step 2: Proof for Brownian bridges.}
Let us next prove the actual statement by comparing the Brownian motions $X^{L,k}$ and $X^{R,k}$ with Brownian bridges $W^{L,k}$ and $W^{R,k}$. We aim at proving that the probability that one of these bridges stays above $-r$ in the first half of the interval when $\Delta_k > (r+W_{G_k})^2 \log^6(1/r)$ is summable when $r$ ranges over $\{2^{-n}, n \in\N\}$, and then that if both $W_{G_k} \leq r \log^{10k}(1/r)$ and $\Delta_k > (r \log^{10k}(1/r))^2$, then there is a probability bounded away from $0$ uniformly in $r$ and $k$ that the two bridges both reach $-r$.
Let us only consider $X^{L,k}$ and $W^{L,k}$ since the same argument applies to $X^{R,k}$ and $W^{R,k}$.
Since, conditionally on $W_{G_k}$, $W_{D_k}$, $\Delta_k$, and $\Delta^L_k$, the path $W^{L,k}$ is a Brownian bridge from $W_{G_k}$ to $W_{D_k}$ with duration $\Delta^L_k \geq \Delta_k$, then the following absolute continuity relation holds: for any measurable and nonnegative function $F$,
\begin{align*}
&\E\bigl[F(W^{L,k}_{t}, t \leq \Delta^L_k/2) \mid W_{G_k}, W_{D_k}, \Delta_k, \Delta^L_k\bigr] \\
& \qquad = \E\Bigl[F(X^{L,k}_t, t \leq \Delta^L_k/2) \, \frac{p_{\Delta^L_k/2}(W_{D_k}-X^L_{\Delta^L_k/2})}{p_{\Delta^L_k}(W_{D_k}-W_{G_k})} \bigm| W_{G_k}, W_{D_k}, \Delta_k, \Delta^L_k\Bigr]
,
\end{align*}
where $p_t(x) = (2\pi t)^{-1/2} \exp(-x^2/(2t))$.
We then integrate with respect to $W_{D_k}$. Precisely, we claim that there exists a constant $C>0$ such that
on the event $\{W_{G_k}>-r, W_{D_k}>-r, \Delta_k > 2(r+W_{G_k})^2 \log^6(1/r)\}$, we have
\begin{equation}\label{eq:borne_ratio_technique}
\E\Bigl[\frac{p_{\Delta^L_k/2}(W_{D_k}-X^L_{\Delta^L_k/2})}{p_{\Delta^L_k}(W_{D_k}-W_{G_k})} \bigm| X^L_{\Delta^L_k/2}, W_{G_k}, \Delta_k, \Delta^L_k\Bigr]
\leq C
.\end{equation}
This allows to infer from the first part of the proof that
\[\P\Bigl(\Delta_k > 2(r+W_{G_k})^2 \log^6(1/r) \enskip\text{and}\enskip \max\Bigl\{\inf_{t \in [0,\Delta_k/2]} W^{L,k}_t, \inf_{t \in [0,\Delta_k/2]} W^{R,k}_t\Bigr\} > -r \Bigr)
\]
is $ O(\log^{-3}(1/r))$. 
Then, as in the first part, we must impose $\Delta_k \leq 2(r+W_{G_k})^2 \log^6(1/r) \leq 3(r \log^{10k+6}(1/r))^2$ for every $k \leq K_r$ in order to prevent the bridges $W^{L,k}$ and $W^{R,k}$ to reach $-r$.

Let us prove the upper bound~\eqref{eq:borne_ratio_technique}. Recall the notation $r_{k}=r \log^{10k}(1/r)$.
Conditionally on $W_{G_k}$ and $\Delta_k$, the path $(W_t; t \in [G_k, D_k])$ has the law of a Brownian motion $Y$ with duration $\Delta_k$, started from $W_{G_k}$, and conditioned on the event $E_k = \{\inf\{Y_t ; t \in [G_k, D_k]\} \leq -r_{k}\}$. Then the law of $W_{D_k}$ can be obtained by splitting at the first hitting time of this value; by the Markov, property, the remaining path after is an unconditioned Brownian motion. Recall also that we work implicitly on the event where $W_{G_k}, W_{D_k} > -r$. 
For $z>0$, let us write $q_z$ for the density of the first hitting time of $-z$ by a standard Brownian motion, 
then:




\begin{align*}
&\E\Bigl[\frac{p_{\Delta^L_k/2}(W_{D_k}-X^L_{\Delta^L_k/2})}{p_{\Delta^L_k}(W_{D_k}-W_{G_k})} \ind{W_{G_k}, W_{D_k}>-r} \bigm| X^L_{\Delta^L_k/2}, W_{G_k}, \Delta_k, \Delta^L_k\Bigr]
\\
&= \frac{1}{\P(E_k)} \int_0^{\Delta_k} \d t \int_{-r}^\infty \d y\ q_{r_{k}+W_{G_k}}(t) p_{\Delta_k-t}(y+r_{k})  \frac{p_{\Delta^L_k/2}(y-X^L_{\Delta^L_k/2})}{p_{\Delta^L_k}(y-W_{G_k})} \ind{W_{G_k}>-r}
\\
&= \frac{1}{\P(E_k)} \int_0^{\Delta_k} \d t \frac{q_{r_{k}+W_{G_k}}(t)}{2\sqrt{\pi(\Delta_k-t)}}  \\
& \qquad  \qquad \qquad \cdot \int_{-r}^\infty \d y \exp\Bigl(\frac{(y-W_{G_k})^2}{2\Delta^L_k} - \frac{(y-X^L_{\Delta^L_k/2})^2}{\Delta^L_k} - \frac{(y+r_{k})^2}{2(\Delta_k-t)}\Bigr) \ind{W_{G_k}>-r}
.\end{align*}
We have $\P(E_k) = \int_0^{\Delta_k} \d t q_{r_{k}+W_{G_k}}(t)$, it thus remains to prove that the integral in $y$ of the rest of the expression is bounded, when in addition $\Delta_k > 2(r+W_{G_k})^2 \log^6(1/r)$.
Let us first consider the integral over $[W_{G_k}, \infty)$. Indeed, if $y > W_{G_k} > -r$, then $(y-W_{G_k})^2 \leq y^2 \leq (y+r_{k})^2$ and consequently
\begin{align*}
&\frac{1}{2\sqrt{\pi(\Delta_k-t)}} \exp\Bigl(\frac{(y-W_{G_k})^2}{2\Delta^L_k} - \frac{(y-X^L_{\Delta^L_k/2})^2}{\Delta^L_k} - \frac{(y+r_{k})^2}{2(\Delta_k-t)}\Bigr)
\\
&\leq \frac{1}{2\sqrt{\pi(\Delta_k-t)}} \exp\Bigl(- \frac{(y+r_{k})^2}{2} \Bigl(\frac{1}{\Delta_k-t}-\frac{1}{\Delta^L_k}\Bigr) - \frac{(y-X^L_{\Delta^L_k/2})^2}{\Delta^L_k}\Bigr)
.\end{align*}
Suppose first that $\Delta^L_k \geq 2 \Delta_k$ so $\Delta^L_k \geq 2 (\Delta_k-t)$ for any $t \in (0, \Delta_k)$ and thus $\frac{1}{\Delta_k-t}-\frac{1}{\Delta^L_k} \geq \frac{1}{2(\Delta_k-t)}$. In this case, by simply removing the last term in the exponential, we have:
\begin{align*}
&\frac{1}{2\sqrt{\pi(\Delta_k-t)}} \exp\Bigl(- \frac{(y+r_{k})^2}{2} \Bigl(\frac{1}{\Delta_k-t}-\frac{1}{\Delta^L_k}\Bigr) - \frac{(y-X^L_{\Delta^L_k/2})^2}{\Delta^L_k}\Bigr)
\\
&\leq \frac{1}{2\sqrt{\pi(\Delta_k-t)}} \exp\Bigl(- \frac{(y+r_{k})^2}{4(\Delta_k-t)}\Bigr)
,\end{align*}
which is the density of a Gaussian law, in particular its integral over the entire real line equals $1$. If $\Delta^L_k$ is close to $\Delta_k$ and $t$ is small, one has to be more careful.
Suppose now that $\Delta_k \leq \Delta^L_k \leq 2 \Delta_k$, then,
\begin{align*}
&\leq \frac{1}{2\sqrt{\pi(\Delta_k-t)}} \exp\Bigl(- \frac{(y+r_{k})^2}{2} \Bigl(\frac{1}{\Delta_k-t}-\frac{1}{\Delta^L_k}\Bigr) - \frac{(y-X^L_{\Delta^L_k/2})^2}{\Delta^L_k}\Bigr)
\\
&\leq \frac{1}{2\sqrt{\pi(\Delta_k-t)}} \exp\Bigl(- \frac{(y+r_{k})^2}{2} \frac{t}{\Delta_k(\Delta_k-t)} - \frac{(y-X^L_{\Delta^L_k/2})^2}{2\Delta_k}\Bigr)
\\
&\leq \frac{1}{2\sqrt{\pi \Delta_k / 2}} \exp\Bigl(- \frac{(y-X^L_{\Delta^L_k/2})^2}{2 \Delta_k}\Bigr) \ind{0 < t \leq \Delta_k/2}\\
& \qquad \qquad \qquad \qquad \qquad  \qquad + \frac{1}{2\sqrt{\pi(\Delta_k-t)}} \exp\Bigl(- \frac{(y+r_{k})^2}{4 (\Delta_k-t)}\Bigr) \ind{\Delta_k/2 < t < \Delta_k}
.\end{align*}
Again the integral over the entire real line of the upper bound is constant.
Finally, if $-r < y < W_{G_k}$ then $(y-W_{G_k})^2 \leq (r+W_{G_k})^2$. Here we use that we work on the event $\Delta_k > 2(r+W_{G_k})^2 \log^6(1/r) \geq 2 (y-W_{G_k})^2 \log^6(1/r)$. Indeed, this implies:
\begin{align*}
&\frac{1}{2\sqrt{\pi(\Delta_k-t)}} \exp\Bigl(\frac{(y-W_{G_k})^2}{2\Delta^L_k} - \frac{(y-X^L_{\Delta^L_k/2})^2}{\Delta^L_k} - \frac{(y+r_{k})^2}{2(\Delta_k-t)}\Bigr)
\\
&\leq \frac{1}{2\sqrt{\pi(\Delta_k-t)}} \exp\Bigl(\frac{1}{4\log^6(1/r)} - \frac{(y+r_{k})^2}{2(\Delta_k-t)}\Bigr)
,\end{align*}
and again the integral over $y$ of the last line is uniformly bounded above.
This completes the proof of the bound~\eqref{eq:borne_ratio_technique}.


Let us finally lower bound the probability that the Brownian bridge $W^{L,k}$ reaches $-r$ in the first half of the interval when both $-r < W_{G_k} \leq r_{k}$ and $\Delta_k > r_{k}^2$, where we keep the notation $r_{k}=r \log^{10k}(1/r)$.
We argue similarly: conditionally on $\Delta^L_k$, $W_{G_k}$, and $W_{D_k}$, this probability equals
\begin{align*}
&\int_0^{\Delta^L_k/2} \d t\ q_{-(r+W_{G_k})}(t) \frac{p_{\Delta^L_k-t}(W_{D_k}+r)}{p_{\Delta^L_k}(W_{D_k}-W_{G_k})}\\
&= \int_0^{\Delta^L_k/2} \d t\ q_{-(r+W_{G_k})}(t) \sqrt{\frac{\Delta^L_k}{\Delta^L_k-t}} \exp\Bigl(\frac{(W_{D_k}-W_{G_k})^2}{2 \Delta^L_k} - \frac{(W_{D_k}+r)^2}{2(\Delta^L_k-t)}\Bigr)
\\
&\geq \int_0^{\Delta_k/2} \d t\ q_{-(r+W_{G_k})}(t) \sqrt{2} \exp\Bigl(- \frac{(W_{D_k}+r)^2}{2 (\Delta_k-t)}\Bigr)
.\end{align*}
Without the exponential, the last integral is bounded away from $0$.
We then lower bound the expectation of the exponential with respect to $W_{D_k}$.
Recall that conditionally on $W_{G_k}$ and $\Delta_k$, the path $(W_t; t \in [G_k, D_k])$ has the law of a Brownian motion started from $W_{G_k}$ and conditioned on the event $E_k$ that it reaches $-r_{k}$ at some point in the time interval, so by splitting at the first hitting time of this value, we have since we restrict our attention to $t \leq \Delta_k/2$,
\begin{align*}
&\E\Bigl[\exp\Bigl(- \frac{(W_{D_k}+r)^2}{2(\Delta_k-t)}\Bigr) \ind{W_{G_k}, W_{D_k}>-r} \bigm| W_{G_k}, \Delta_k\Bigr]
\\
&= \frac{1}{\P(E_k)} \int_0^{\Delta_k} \d s \int_{-r}^\infty \d y\ \frac{q_{-(r_{k}+W_{G_k})}(s)}{\sqrt{2\pi(\Delta_k-s)}} \exp\Bigl(- \frac{(y+r_{k})^2}{2(\Delta_k-s)} - \frac{(y+r)^2}{2 (\Delta_k-t)}\Bigr) \ind{W_{G_k}>-r}
\\
&\geq \int_0^{\Delta_k} \d s \int_{-r}^\infty \d y\ \frac{q_{-(r_{k}+W_{G_k})}(s)}{\sqrt{2 \pi(\Delta_k-s)}} \exp\Bigl(- \frac{(y+r_{k})^2}{2} \Bigl(\frac{1}{\Delta_k-s}+\frac{1}{\Delta_k-t}\Bigr)\Bigr) \ind{W_{G_k}>-r}
\\
&\geq \int_0^{\Delta_k/2} \d s \int_{-r}^\infty \d y\ \frac{q_{-(r_{k}+W_{G_k})}(s)}{\sqrt{2 \pi \Delta_k}} \exp\Bigl(- 2 \frac{(y+r_{k})^2}{\Delta_k}\Bigr) \ind{W_{G_k}>-r}
.\end{align*}
Consider the integral over $y$ only in the last line, and without de density $q$. It equals half the probability that a Gaussian random variable with mean $-r_{k}$ and variance $\Delta_k/4$ is larger than $-r$, or equivalently that a Gaussian random variable with mean $0$ and variance $1$ is larger than $2(r_{k}-r)/\sqrt{\Delta_k}$. Since $\Delta_k > r_{k}^2$, this probability is indeed bounded away from $0$ uniformly in $k$ and $r$.
Then the integral of $q$ over $s$ equals the probability that a standard Brownian motion reaches $-(W_{G_k}+r_{k})$ in less than $\Delta_k/2$ when $W_{G_k} \leq r_{k} \leq \sqrt{\Delta_k}$, which is again bounded away from $0$.

Combined with the results on $X^{L,k}$, we infer that there exists some constant $c > 0$ such that for all $r<1$, on each interval $[G_k,D_k]$ with $W_{G_k} \leq r_{k}$ and $\Delta_k \geq (r_{k})^2$, there is a probability at least $c$ that $W^{L,k}$ reaches the value $-r$ in the first half of the interval. The same holds for $X^{L,k}$ and since they are independent, they both reach the value $-r$ in the first half with probability at least $c^2$. We can then conclude from Lemma~\ref{lem:ruine} that this occurs for a positive proportion of indices $k \leq K_r$.
\end{proof}

\end{document}
